% ===========================================================================
\ReportSection[sec:hw]{Hardware, Price Points and Configurations}

\ReportSubsection[sec:nas]{Computer 1: NAS and private-cloud server}

\TableTitle{Computer 1 reference bill of materials (planning estimates)}
{\footnotesize
\begin{xltabular}{\textwidth}{@{}Y{0.621}Y{1.379}r@{}}
  \toprule
  Component & Specification & Est. cost \\
  \midrule
  \endhead
  Chassis and platform & 2U/4U rack, 12 or more 3.5-inch bays, one 16--24-core server CPU, TPM~2.0, BMC, redundant power supplies & \$3--4K \\
  Memory & 256~GB ECC & \$1.5--2.5K \\
  System and VM storage & 2 $\times$ 1.9~TB NVMe (mirror, OS and hypervisor); 2 $\times$ 3.8~TB NVMe (mirror, VMs, vector database) & \$1--1.5K \\
  Bulk storage & 8 $\times$ 20~TB enterprise HDD in RAID-Z2 or RAID~6 (about 100~TiB usable; plan at most 80\% fill); 4 bays free for growth & \$3--3.6K \\
  Controller and network & SAS HBA; dual-port 25~GbE & \$0.5--0.8K \\
  UPS and rack & 3~kVA UPS shared with the ML nodes & \$1--2K \\
  OS and hypervisor subscription & FIPS-capable Linux and virtualization support & \$0.3--1K \\
  Switch, firewall, cabling & 25~GbE switch with spare ports, firewall, management network & \$1.5--2.5K \\
  \midrule
  \textbf{Computer 1 total} & & \textbf{\$12--18K (about \$15K)} \\
  \bottomrule
\end{xltabular}}

\textbf{Storage sizing (best estimate, to confirm).} The archive is assumed to be about
20--40~TB now, growing 20--25\% per year, so about 40--70~TB usable is needed within 3
years including snapshots, vector indexes and model artifacts. The 8 $\times$ 20~TB array
(about 100~TiB usable, filled to at most 80\%) covers this with room to spare, and 4
bays remain free.

\TodoAnchor{aoa-backup}
\textbf{Redundancy.} Following the sponsor's direction, Computer~1 is protected by its
disks rather than by a second server: RAID-Z2 (tolerates two disk failures), a hot spare
(add one 20~TB drive, about \$0.4K), redundant power supplies and a UPS, plus snapshots.
This covers disk failure and accidental deletion. It does not cover loss of the server,
ransomware or a site loss, because the data and the snapshots share one machine; an
encrypted offline backup set rotated off-site (about \$2--4K) closes that gap and is
recommended but optional.

Likely additions not included in the table: Windows Server, database and client-access
licensing for Autodesk and Bluebeam server components (about \$2--8K; estimate, depends on
the products chosen; Section~\ref{sec:licenses}).

\ReportSubsection{Computer 2: ML node building blocks}

\TableTitle{ML node options (each in a 4-GPU-capable chassis so GPUs can be added later)}
{\footnotesize
\begin{xltabular}{\textwidth}{@{}l Y{0.833} Y{1.167} r@{}}
  \toprule
  Node & GPUs & Host and notes & Est. cost \\
  \midrule
  \endhead
  \textbf{N1} & 1 $\times$ RTX PRO 6000 Blackwell 96~GB (Max-Q, 300~W), \$8--10K & 16--32-core CPU, 256~GB RAM, 2 $\times$ 3.8~TB NVMe, dual 1600--2000~W PSUs, 25~GbE: \$9--11K; misc.\ \$1K & \$18--22K \\
  \textbf{N2} & 2 GPUs (192~GB total), \$16--20K & Same class of host: \$10--11K; misc.\ \$1K & \$27--32K \\
  \textbf{N4} & 4 GPUs (384~GB total), \$32--40K & Host with 512~GB RAM, larger PSUs: \$12--13K; misc.\ \$2K & \$46--55K \\
  \bottomrule
\end{xltabular}}

Adding one GPU to an existing node costs about \$8--10K. Government-channel
(TAA-compliant OEM) pricing is often 15--40\% higher than the street prices used
here (Section~\ref{sec:procure}).

\ReportSubsection{System configurations and price ladder}

\TableTitle{Configurations (Computer~1 plus ML nodes)}
{\footnotesize
\begin{xltabular}{\textwidth}{@{}l Y{0.8} r r Y{1.2}@{}}
  \toprule
  Config. & ML node(s) & Hardware total & GPU memory & Named users \\
  \midrule
  \endhead
  \textbf{S1} & N1 & \$35K (\$30--40K) & 96~GB & 10--30 \\
  \textbf{S2} & N2 & \$44K (\$39--50K) & 192~GB & 25--75 \\
  \textbf{D1} & N1 + N1 (A + B) & \$55K (\$48--62K) & 2 $\times$ 96~GB & 10--30 \\
  \textbf{D3} & N2 (A) + N1 (B) & \$64K (\$57--72K) & 192 + 96~GB & 30--100 \\
  \textbf{S4} & N4 & \$65K (\$58--73K) & 384~GB & 75--250 \\
  \textbf{D4} & N2 + N2 (A + B) & \$73K (\$66--82K) & 2 $\times$ 192~GB & 30--100 (A); up to 60--200 with both serving \\
  \bottomrule
\end{xltabular}}

\TableTitle{Upgrade paths}
{\small
\begin{tabularx}{\textwidth}{@{}Y{1.0}Y{0.6}Y{1.4}@{}}
  \toprule
  Step & Added cost & Result \\
  \midrule
  S1 $\to$ S2 (add one GPU) & about \$9K & Doubles memory; train and serve at once on separate GPUs. \\
  S1 $\to$ D1 (add node B) & about \$20K & Dedicated training node and failover path. \\
  D1 $\to$ D3 (add GPU to node A) & about \$9K & Roughly doubles serving capacity; failover is at reduced capacity. \\
  D3 $\to$ D4 (add GPU to node B) & about \$9K & Failover at full capacity; both nodes can serve when not training. \\
  Any $\to$ third node & \$18--50K & More replicas or parallel training jobs. \\
  \bottomrule
\end{tabularx}}

\ReportSubsection[sec:dell]{Dell-configured pricing}
\TodoAnchor{aoa-dell}

Hardware would be ordered from Dell, configured to specification. Street prices
understate that cost because enterprise servers include vendor-qualified components,
warranty and support. A planning multiplier of 1.3 (range 1.15--1.4) is applied to
hardware until a quote is available; labor and operating costs are unchanged.

\TableTitle[tbl:dell]{Hardware and 3-year cost with Dell-configured pricing (USD)}
{\footnotesize
\begin{xltabular}{\textwidth}{@{}l Y{0.848} Y{1.091} Y{0.848} Y{1.213}@{}}
  \toprule
  Config. & Street-price hardware & Dell-configured hardware (range) & Dell hardware (expected) & 3-yr total, Dell-configured \\
  \midrule
  \endhead
  S1 & \$30--40K & \$35--56K & \$46K & \$81--121K \\
  S2 & \$39--50K & \$45--70K & \$57K & \$94--140K \\
  D1 & \$48--62K & \$55--87K & \$72K & \$117--174K \\
  \textbf{D3} & \$57--72K & \$66--101K & \textbf{\$83K} & \textbf{\$131--192K} \\
  S4 & \$58--73K & \$67--102K & \$85K & \$122--181K \\
  D4 & \$66--82K & \$76--115K & \$95K & \$144--211K \\
  \bottomrule
\end{xltabular}}
{\footnotesize 3-year totals add 3 years of operating cost, setup labor and upkeep (Table~\ref{tbl:tco}) to the Dell-configured hardware.}

\textbf{Notes for the Dell quote.}
\begin{itemize}
  \item Ask which current GPU-capable PowerEdge models (for example the 2U R-series and the
    GPU-dense XE-series) and Precision rack or tower workstations support the RTX PRO 6000
    Blackwell, how many GPUs each holds, and the power and cooling needed. The GPU count
    per chassis decides the growth path; confirm model numbers with the Dell account team.
  \item Dell servers typically use passive 600~W Server Edition GPUs instead of 300~W Max-Q
    cards; add about 0.3~kW per GPU to the power figures in the facility table.
  \item For Computer~1, request a storage-dense 2U PowerEdge with 12 or more 3.5-inch bays,
    redundant power, TPM~2.0 and iDRAC, with NVMe boot or VM drives.
  \item Dell support bundled in the quote (for example ProSupport) may replace part of the
    annual non-labor support cost assumed below.
  \item Request country-of-origin and TAA/Section~889 documentation with the quote.
\end{itemize}

\ReportSubsection[sec:sizing]{Sizing logic}
\TodoAnchor{aoa-bench}

\textbf{Weights memory} is parameters times bytes per parameter (FP16 = 2,
FP8 = 1, 4-bit $\approx$ 0.5), plus about 10\% overhead. The \textbf{KV cache}
per token is
\[
  \text{KV bytes/token} = 2 \times \text{layers} \times \text{KV heads} \times \text{head dim} \times \text{bytes per element}.
\]

\textbf{Example 1: gpt-oss-120b on one 96~GB GPU.} Approximate architecture: 36
layers, 8 KV heads, head dimension 64, with about half the layers using a short
sliding window, so only about 18 layers cache the full context. In BF16 the
cache is about $2 \times 18 \times 8 \times 64 \times 2~\text{B} \approx 37$~KB per
token, or about 4.7~GB for a 128K-token session. With
$96 \times 0.9 \approx 86$~GB and about 63~GB of weights, roughly 23~GB remains
for cache: about 5 simultaneous 128K sessions or 70 or more at 8K context.
Compute, not memory, usually limits short-prompt use, which is why one GPU is
rated for 3--8 simultaneous active users.

\textbf{Example 2: 70B dense model in FP8 on two GPUs (192~GB).} With 80 layers,
8 KV heads, head dimension 128 and an FP8 cache, the cost is
$2 \times 80 \times 8 \times 128 \times 1~\text{B} \approx 160$~KB per token (about
21~GB per 128K session). After about 70~GB of weights, about 103~GB remains:
roughly 5 full-128K sessions or 80 at 8K. A mixture-of-experts model is therefore
much more efficient than a dense 70B model for long-context, multi-user serving.

\[
  \text{Concurrent sessions} \approx \frac{\text{GPU memory} \times 0.9 - \text{weights}}{\text{average context tokens} \times \text{KV bytes per token}}
\]

\textbf{Throughput (order of magnitude).} Each GPU has about 1.8~TB/s memory
bandwidth. A small-active-parameter mixture-of-experts model can sustain roughly
1,000--3,000 aggregate output tokens/s per GPU under batching. Dense 70B models
are several times slower. \emph{Benchmark the actual model and real prompts
before purchase.}

\textbf{Fine-tuning memory (planning estimates).}

\TableTitle{Fine-tuning feasibility on 96 GB and 192 GB nodes}
{\small
\begin{xltabular}{\textwidth}{@{}Y{1.138}Y{0.828}Y{1.034}@{}}
  \toprule
  Method and model & Approx.\ GPU memory & Fits on \\
  \midrule
  \endhead
  Adapter (LoRA/QLoRA), 7--30B model & 20--70~GB & One 96~GB GPU \\
  QLoRA, 70B dense model (4-bit base) & 48--70~GB & One 96~GB GPU (tight at long sequence lengths) \\
  Adapter fine-tuning of gpt-oss-120b or similar & about 65--90~GB or more (verify with current tooling) & One 96~GB GPU is tight; two GPUs (192~GB) is comfortable \\
  Full fine-tuning (about 16 bytes per parameter with Adam in mixed precision) & 7B: about 112~GB; 30B: about 480~GB & 7B or smaller only on a 2-GPU node; larger is out of scope \\
  Pre-training a new LLM & Far beyond this budget & Not feasible \\
  \bottomrule
\end{xltabular}}

Adapter jobs on a few thousand to roughly one hundred thousand examples typically
run for hours to a few days on one GPU; confirm in the pilot. The GPUs connect
over PCIe (no NVLink), which is fine for serving and adapter training.

\ReportSubsection{Sizing check against the stated requirement}

Requirement: at least 20 named users, about 12 generating at the same time.
\begin{itemize}
  \item \textbf{Streaming load.} 12 users $\times$ 20 tokens/s is about 240 tokens/s,
    far below the roughly 1,000--3,000 tokens/s a mixture-of-experts model can sustain per
    GPU in aggregate.
  \item \textbf{Cache memory.} With gpt-oss-120b, 12 sessions of 16K tokens need about 7~GB;
    even 3 sessions at 128K (about 14~GB) plus 9 at 16K (about 5~GB) need about 19~GB, which
    fits in the roughly 23~GB left on one 96~GB GPU after the weights.
  \item \textbf{Limiting factor.} Processing many long prompts at once (prompt prefill) and
    keeping time to first token low. This is why one GPU is rated 3--8 simultaneous users;
    the rating is conservative and the pilot benchmark may show more.
  \item \textbf{Conclusion.} A 1-GPU serving node is marginal for 12 users. A 2-GPU node
    (rated 8--20) meets the need with margin, and Node~B can add capacity when it is not
    training. If named users grow well beyond 20, simultaneous use may exceed 20: add a GPU
    or use Node~B as a second serving node.
\end{itemize}

\ReportSubsection{Facility}
\TodoAnchor{aoa-facility}

\TableTitle{System power at load and circuit guidance}
{\small
\begin{tabularx}{\textwidth}{@{}l Y{0.6} Y{1.4}@{}}
  \toprule
  Config. & Power at load & Circuit and cooling \\
  \midrule
  S1 & 1.1--1.5~kW & One 208/240~V 20~A circuit. \\
  S2 & 1.6--2.1~kW & One 208/240~V 20~A circuit. \\
  D1 & 1.9--2.5~kW & One 208/240~V 20~A circuit (about 3.3~kW continuous). \\
  D3 & 2.4--3.1~kW & Two 20~A circuits, or one 30~A circuit, with separate PDUs per node. \\
  S4 & 2.3--3.1~kW & Two 20~A circuits or one 30~A circuit. \\
  D4 & 2.9--3.7~kW & Two 20~A circuits or one 30~A circuit; about 12,600~BTU/h at the top of the range. \\
  \bottomrule
\end{tabularx}}

Figures assume 300~W (Max-Q) GPUs: about 0.8--1.0~kW for N1, 1.3--1.6~kW for N2,
2.0--2.6~kW for N4 and 0.3--0.5~kW for Computer~1. If Dell supplies 600~W Server
Edition GPUs, add about 0.3~kW per GPU. The nodes are loud; place them
in a machine room or sound-attenuated rack.

\ReportSubsection[sec:procure]{Procurement considerations}
\begin{itemize}
  \item \textbf{Dell as vendor; TAA, Section 889 and NDAA compliance.} Dell is the stated
    vendor for configured-to-order servers. Request quotes for Computer~1 and for each
    candidate node configuration, with country-of-origin and TAA/Section~889
    documentation, and compare with the street-price planning figures in this report
    (Section~\ref{sec:dell}).
  \item \textbf{Memory and disk price volatility.} Server RAM and large HDD prices
    have been unstable. Quote early and consider buying less RAM now with expansion
    later.
  \item \textbf{Warranty and spares.} Three-year on-site coverage where available; hold
    a spare PSU and a spare HDD at minimum. Plan to retain failed drives (CUI)
    rather than returning them.
  \item \textbf{License servers.} Confirm with each software vendor that the license
    server may run in a virtual machine and what hardware binding applies, before
    finalizing the Computer~1 design.
  \item \textbf{Lead time and obsolescence.} Confirm GPU availability; plan for a
    3--5 year service life. The model layer can be replaced without hardware change.
\end{itemize}
